\title{LongRoPE: Extending LLM Context Window Beyond 2 Million Tokens}

\begin{abstract}
A large context window is a sought-after capability in large language models (LLMs). Nevertheless, existing approaches for expanding context windows beyond approximately 128k tokens face significant barriers, such as the elevated cost of fine-tuning, limited availability of sufficiently lengthy training data, and instability resulting from the introduction of new token positions.

This work presents LongRoPE, marking the inaugural approach to extend the context window of existing pre-trained LLMs to an extraordinary \textbf{2048k} tokens, accomplished with as few as 1k fine-tuning steps at training lengths up to 256k, all while preserving accuracy at the standard short context window. This capability arises from three principal advancements: (i) we pinpoint and leverage two kinds of positional interpolation non-uniformities via an optimized search mechanism, facilitating superior fine-tuning initialization and offering an 8$\times$ expansion without further fine-tuning; (ii) we propose a progressive expansion method, beginning with fine-tuning at 256k length, and subsequently applying a secondary positional interpolation on the already fine-tuned model to realize the full 2048k context window; (iii) we further refine LongRoPE at the 8k length to fully recover baseline short-window performance. Comprehensive evaluations using LLaMA2 and Mistral models on multiple tasks illustrate the robustness and generality of our approach. Models adapted with LongRoPE retain their original structures except for minor updates to positional embeddings, ensuring compatibility with most previously implemented optimizations. The codebase will be released at \texttt{https://github.com/microsoft/LongRoPE}

\section{Introduction}

Although Large Language Models (LLMs) have achieved outstanding performance across a range of tasks~\cite{gpt4,llama2}, they are constrained by their finite context window size, such as the 4096-token limit of LLaMA2~\cite{llama2}. The effectiveness of LLMs diminishes once inputs exceed this context window, as models have not been trained to process these additional positions. This constraint creates difficulties for key applications, including in-context learning with large numbers of examples~\cite{Huang2023FewerIM} and the deployment of LLM agents~\cite{park2023generative,madaan2023selfrefine}.

Recent studies demonstrate that the context window of a pre-trained LLM can be pushed to approximately 128k tokens through targeted fine-tuning with extended text sequences~\cite{longlora,pi,yarn,zhang2024soaring,liu2023scaling}. However, scaling the context window beyond this faces several substantial challenges. First, when introducing untrained positional indices during window enlargement, a significant number of abnormal values emerge, causing out-of-distribution problems and impeding stable convergence during the fine-tuning process~\cite{pi}. This difficulty is heightened in cases where an increase from 4k to over $1000k$ tokens adds more than 90\% entirely novel positions. Second, fine-tuning generally necessitates textual data that matches these extensive lengths. The current corpora have a shortage of extremely lengthy passages—those above 1000k tokens—making it difficult to source suitable data. In addition, handling such exceptionally long texts demands considerable computational resources and training time, placing immense strain on available GPU capacity. Third, at such expansive context windows, attention mechanisms distribute focus very broadly, resulting in diluted attention across many positions; this dispersion ultimately undermines model performance on tasks involving shorter contexts~\cite{pi}.

A common strategy to address the initial challenge involves interpolating RoPE positional embeddings~\cite{rope,pi}, which rescales novel position indices to align them with the pretrained positional interval, as illustrated in Fig.\ref{fig:overview}. The Position Interpolation (PI)~\cite{pi} technique modifies RoPE's rotary angles through linear interpolation according to the extension ratio. The NTK method~\cite{ntk,dynamicntk} introduces the idea of applying non-uniform interpolation and extrapolation along RoPE’s various dimensions. YaRN~\cite{yarn} further extends this by dividing RoPE dimensions into three frequency-driven clusters, employing a combination of extrapolation, NTK, and linear interpolations for each group. Nonetheless, positional embeddings in Transformer models contain a \emph{complex, non-uniform pattern of information entropy}. Current methods do not fully exploit this nuanced non-uniformity, resulting in information degradation that constrains the effective expansion of the context window.

Section~\ref{sec:analysis} presents two principal empirical observations: \textbf{(1)} To achieve effective positional interpolation, it is important to address two types of non-uniformity: discrepancies in RoPE dimension scaling and in the distribution of token positions. For dimensions at the lower end of the RoPE spectrum and early token positions, minimal interpolation is advantageous, yet the precise optimal configuration is contingent upon the intended sequence length extension.  
\textbf{(2)} Incorporating these non-uniformities into the interpolation strategy facilitates the preservation of crucial information from the original RoPE—specifically, certain key dimensions and specific token positions. As a result, the adverse effects stemming from interpolation are mitigated, yielding superior initialization for subsequent fine-tuning. Furthermore, this approach enables sequence lengths to be extended by a factor of 8 in settings without fine-tuning.

Building upon these observations, we propose \textbf{LongRoPE}, a robust approach that expands the LLM context window to surpass 2 \emph{million} tokens. The design of LongRoPE relies on three primary innovations. To begin with, {LongRoPE} takes full advantage of multidimensional, non-uniform features within positional interpolation. It determines optimal rescale factors for the rotation angles in each RoPE dimension according to token positions. Since the process of identifying these rescale factors becomes exponentially more complex as the desired extension ratio increases, {LongRoPE} employs an evolutionary search method enhanced by two optimization strategies to improve the efficiency of the search. An illustrative example of the resulting rescaled RoPE identified via this process is provided in Fig.~\ref{fig:overview}.

Subsequently, {LongRoPE} employs a systematic and efficient progressive extension approach, enabling it to reach a 2048k context window without requiring direct fine-tuning on ultra-long texts, which are rare and difficult to obtain. The approach initiates by exploring a 256k context length using the pre-trained LLM, followed by fine-tuning at this window size. Taking advantage of our non-uniform positional interpolation—which permits an eightfold expansion even without further fine-tuning—we then perform a subsequent search for optimal RoPE rescale factors using the previously fine-tuned, extended LLM. Through this process, the method ultimately supports a 2048k context window for both LLaMA2 and Mistral~\cite{mistral}.

In order to prevent loss of performance on the initial, shorter context windows, {LongRoPE} further modifies the RoPE rescale factors for the expanded LLM. Analogous to the upscaling process from 256k to 2048k, we apply our search algorithm to rescale down to 4k and 8k context window sizes for the 256k fine-tuned LLM, aiming to minimize positional interpolation. When performing inference, if the input sequence length falls below 8k, RoPE is updated with the rescale factors identified through this search.

A comprehensive set of experiments conducted on multiple LLMs and a range of long-context evaluation tasks validate the efficacy of our approach. Our results indicate that LongRoPE consistently sustains low perplexity scores from 4k up to 2048k sequence lengths, secures more than 90\% passkey retrieval accuracy, and achieves similar performance on conventional benchmarks situated within the 4096-token context window. Furthermore, LongRoPE is compatible with any LLM employing RoPE embeddings. We intend to make our code and LongRoPE-2048k models publicly available.

\section{Non-uniformity in Positional Interpolation}
\label{sec:analysis}

\subsection{Preliminary}

Transformer architectures inherently lack positional awareness and, therefore, depend on explicit position embeddings to encode token order within input sequences. In this study, we investigate the RoPE~\cite{rope} positional embedding approach, a method commonly adopted by state-of-the-art LLMs. For a token located at position index $n$, its RoPE embedding is succinctly described as:

\begin{equation}
		\fontsize{7.8}{7.8} \selectfont
	\label{eq:rope}
	\left[ cos(n\theta_0), sin(n\theta_0), cos(n\theta_1), \cdot\cdot\cdot, cos(n\theta_{d/2-1}), sin(n\theta_{d/2-1}) \right]   
\end{equation}

Here, $d$ denotes the embedding dimension, 
 $n\theta_i$ corresponds to the rotational angle assigned to the token at position $n$, and
$\theta_i = \theta^{-2i/d}$ specifies the rotation frequencies utilized. In RoPE, $\theta$ is typically set to a default value of 10000.

\textbf{\emph{Context window extension ratio $s$} and positional interpolation}. We define $s$ as  the ratio of extended context length $L'$ to the original length $L$: $s=\frac{L'}{L}$. 

In order to increase the context window from $L$ to $L'$, existing positional interpolation techniques advocate reducing the rotation frequencies $\theta_i$ in proportion to the extension ratio $s$. Defining $\beta = \theta^{2/d}$ and using $\lambda$ to represent the actual scaling factor associated with $s$, we can present a generalized form for these positional interpolation approaches as:

\begin{equation}	
	\fontsize{7.5}{7.5} \selectfont
	\label{eq:unified_rope}
	\begin{aligned}
		\left[cos\left(\frac{n}{\lambda(\beta)^0}\right), sin \left(\frac{n}{\lambda(\beta)^0}\right), cos\left(\frac{n}{\lambda(\beta)^1}\right), \cdots,  sin \left(\frac{n}{\lambda(\beta)^{d/2-1}}\right) \right]   
	\end{aligned}
\end{equation}

\textbf{Linear positional interpolation (PI)}. PI~\cite{pi} proposes performing a linear interpolation on position indices as long as they remain within the pre-training context window. When extending to a new sequence length by a factor of $s$, this approach scales down the rotation angles for every position uniformly by a factor $\lambda=s$ in every RoPE dimension. Nonetheless, this uniform compression causes positional representations to become overly concentrated, making it challenging for the model to discriminate between tokens that are spatially proximate. Consequently, PI generally yields suboptimal results as the extension ratio increases.

\textbf{NTK-based interpolation and extrapolation}. ~\cite{ntk,dynamicntk} investigate RoPE from the standpoint of information encoding, employing Neural Tangent Kernel (NTK) theory~\cite{jacot2018neural,tancik2020fourier} in their analysis. To address the problem of congested positions in PI, they propose redistributing the interpolation demand over the different dimensions of RoPE. Specifically, they attenuate the scaling in lower (high frequency) dimensions while amplifying it in higher (low frequency) dimensions. This approach facilitates both interpolation and extrapolation over positions, expressed as $\lambda=s^{i}$. In an advancement, the dynamic NTK~\cite{dynamicntk} variant modulates the extension ratio at each position relative to the observed sequence length. Distinct from PI, which typically requires fine-tuning, NTK-informed strategies permit the extension of context lengths without additional training; however, their extension ratios generally reach up to 4$\times$.

\textbf{YaRN}~\cite{yarn} offers a notable enhancement to positional interpolation by categorizing RoPE dimensions into three distinct frequency ranges, with each group subjected to a specific interpolation method. Specifically, extrapolation ($\lambda$=1) is applied to dimensions with the highest frequencies, linear interpolation (PI) is used for the lowest frequency dimensions, and those with intermediate frequencies utilize NTK. Central to YaRN’s approach is its dimension grouping methodology, which at present is based on manually crafted empirical heuristics. This manual process could limit its adaptability and potentially hinder optimal performance when deploying YaRN to novel LLM architectures.

\subsection{Study on Non-uniform Positional Interpolation}

Building on insights from NTK and YaRN, we observe that their performance improvements stem from employing non-linear approaches—particularly by attending to distinct frequencies in RoPE dimensions for tailored interpolation and extrapolation. Yet, prevailing non-linear methods are predominantly constructed based on heuristics. This brings forth two pertinent questions: (1) Does the present form of positional interpolation represent an optimal solution? (2) Are there additional, yet-untapped, forms of non-linearity?

In addressing these questions, we employ evolution search (refer to Sec.~\ref{sec:method}) to identify improved non-uniform positional interpolations for LLaMA2-7B. This optimization process is steered by evaluating perplexity, utilizing five randomly selected instances from the PG19~\cite{pg19} validation dataset. Our experiments yield several important observations.

\textbf{Finding 1}: \textit{RoPE dimensions exhibit substantial non-uniformities, which are not effectively handled by current positional interpolation methods.}

We explore the optimal value of $\lambda$ for each RoPE dimension as presented in Eq.~\ref{eq:unified_rope}. Table~\ref{tbl:exp1} reports the perplexity of LLaMA2-7B, evaluated on PG19 and Proof-pile~\cite{proof-pile} test sets, using several methods without any fine-tuning. Our approach yields notable improvements over existing methods, indicating that both current linear (PI) and non-uniform (Dynamic-NTK and YaRN) interpolation strategies fail to achieve optimality. Interestingly, on PG19, YaRN delivers worse performance compared to PI and NTK, primarily because it cannot attain the target context window length for LLMs that are not fine-tuned; specifically, for an 8k context, YaRN’s perplexity increases sharply after 7k tokens.

In our investigation, we observe that the rescaling factors $\lambda$ in Eq.~\ref{eq:unified_rope} are not uniform, unlike the constant scale $s$ utilized in PI, the formulaic approach in NTK, or the group-based method in YaRN. This non-uniformity in scaling leads to marked improvements in LLaMA2’s language modeling capability—specifically, perplexity—when handling 8k and 16k token context windows, even in the absence of fine-tuning. The underlying reason is that these adjusted positional embeddings maintain much of the structure inherent to the original RoPE, particularly in key dimensions, thereby making it less challenging for large language models to differentiate between tokens with proximate positions.

\textbf{Finding 2}: \textit{RoPE for the initial tokens in the input sequence should be extrapolated with less interpolation.} 

We propose that the RoPE interpolation should be reduced for the first $\hat{n}$ tokens in input sequences, as these tokens tend to draw large attention weights and play a vital role in the functioning of attention mechanisms—a trend also noted in Streaming LLM~\cite{streamingllm} and LM-Infinite~\cite{han2023lminfinite}. To assess this hypothesis, we perform experiments by expanding the context window to 8k and 16k with both PI and NTK approaches, while deliberately excluding the initial $\hat n$ tokens (for $\hat n=0,2,...,256$) from interpolation. The scenario with $\hat n=0$ is equivalent to conventional PI and NTK methods. Our results, summarized in Table~\ref{tbl:tokenposition}, reveal two principal findings: \textbf{(1)} excluding interpolation for the leading tokens in the sequence provides a tangible performance benefit, and \textbf{(2)} the optimal value of $\hat n$ is sensitive to the specific target context length.

\textbf{Finding 3}: \textit{Non-uniform positional interpolation effectively extends LLM context window in both fine-tuning and non-fine-tuning settings.} 

Our investigation demonstrates that employing our optimized non-uniform position interpolation leads to substantial gains in context extension to 8k and 16k lengths without additional fine-tuning. However, achieving further increases in context window size necessitates a fine-tuning step. Accordingly, we perform fine-tuning of LLaMA2-7B using our tailored RoPE strategy for a 64k token context window (see Appendix for experimental details). As indicated in Table~\ref{tbl:finetune}, our approach consistently surpasses both PI and YaRN methods, not only in the pre-fine-tuning stage but also following fine-tuning of LLaMA2-7B. This performance advantage results from the strategic implementation of non-uniform positional interpolation, which reduces information loss and offers a superior starting point for subsequent fine-tuning.

\textbf{Summary}. This work identifies two key sources of non-uniformity: variation among RoPE dimensions and differences among token positions. Properly leveraging these non-uniform factors in positional interpolation markedly boosts LLM performance for extended context windows.

\section{LongRoPE}

\label{sec:method}
Building upon these insights, we propose LongRoPE, which initially implements an effective search algorithm designed to leverage both types of non-uniformities. This approach subsequently enables the extension of the LLM context window to more than 2 million tokens.

\subsection{Problem Formulation}

The presence of these two types of non-uniformity significantly expands the solution space and adds layers of difficulty to the optimization process. To overcome these challenges, we recast the problem of optimizing multidimensional non-uniform positional interpolation as a search problem.

Consider a LLM that is designed to handle a context window of length $L'$ and is provided with input documents $\mathbf{X}$, in which each sequence $\mathbf{x}\in\mathbf{X}$ exceeds $L'$ tokens. For such cases, we represent the native rotary angle for the $i^{th}$ component in the RoPE embedding at a given token position $n$ as $\frac{n}{\beta^i}$. Based on this, the optimization task is articulated as follows:

\begin{equation}
\begin{gathered}
		\label{eq:problem}

\underset {\mathbf{x}\in\mathbf{X}; \, |\mathbf{x}|\ge L'}{ \operatorname {arg\,min} } \,  \mathcal{L}\, \left( \text{LLM} (\text{RoPE}, \mathbf{X})\right),	\text{where \;} 
\\
\scalebox{0.93}{
$\underset {\begin{subarray}{l} i=0,\ldots,\frac{d}{2}-1;\\ n\in[ 0,|\mathbf{x}|);\end{subarray}}{\text{RoPE}_(n)}= {\left[\ldots, \cos\left(\mathbb{I}(\hat{\lambda}_{i}, \hat{n})\frac{n}{\beta^i}\right), \sin\left(\mathbb{I}(\hat{\lambda}_{i}, \hat{n})\frac{n}{\beta^i}\right),\ldots\right] }$}\\
	\text{with \;} \mathbb{I}(\hat{\lambda}_{i},\hat{n})=\begin{cases}
	1&\mbox{\quad  $n<\hat{n}$}\\
{\frac{1}{\lambda}_{i}}&\mbox{\quad $n\geq\hat{n}$}
\end{cases}
	\end{gathered}
\end{equation}

In this approach, we define a set of rescaling factors, $\mathbb{I}(\hat{\lambda}_{i},\hat{n})$, designed to address both types of non-uniformities. Here, $\hat{\lambda_i}$ corresponds to non-uniformities present in RoPE dimensions, while $\hat{n}$ represents those related to token positions. The rescaling function $\mathbb{I}(\hat{\lambda}_{i},\hat{n})$ adjusts the rotation angle for the $i^{th}$ RoPE dimension; in this context, $\hat{\lambda}_{i}$ acts as the scaling parameter, and $\hat{n}$ specifies the position threshold for tokens. For the initial positions, that is, the first $\hat{n} - 1$ tokens, $\hat{\lambda}_{i}$ does not alter the computation, so the original RoPE rotary angle $\frac{n}{\beta^i}$ is retained. Once token positions satisfy $n \ge \hat{n}$, the rescaling factor becomes active.

With a specified target context window length of $L'$, our aim is to determine the set of optimal rescaling coefficients ($\mathbb{I}(\hat{\lambda}_{0},\hat{n})$, $\mathbb{I}(\hat{\lambda}_{1},\hat{n})$, ...$\mathbb{I}(\hat{\lambda}_{i},\hat{n})$, ...) across the $1^{st}$ through $d^{th}$ dimensions of RoPE. The intention is that, when applying the rescaled $\text{RoPE}$ to the target $\text{LLM}$, the resulting model minimizes the next token prediction loss, $\mathcal{L}$ (interpreted here as perplexity), over input datasets $\mathbf{X}$ that have a token sequence length of $L'$.

\subsection{Searching the Non-uniform Position Interpolation}

To address the problem formulated in Eq.~\ref{eq:problem}, we present a straightforward yet powerful approach that identifies the optimal RoPE rescale factors, thereby leveraging the multidimensional variability inherent in position embedding.

\textbf{Search space}. We construct an extensive search space that accommodates both types of non-uniformity, as depicted in Table~\ref{tbl:searchspace}. In particular, our method permits the independent tuning of a rescale factor for each dimension within the RoPE embedding. For the sake of simplicity in search space construction, we opt to search over $\lambda_i$ and $\hat{n}$, rather than directly searching for $\mathbb{I}(\hat{\lambda}_{i},\hat{n})$, given that $\hat{\lambda}_{i}=1/\lambda_i$. According to Table~\ref{tbl:searchspace}, the possible values for $\lambda_i$ range from a lower bound of 1.0 (which equates to direct extrapolation) up to an upper bound of $s\times1.25$ (allowing greater interpolation compared to PI), with increments of 0.01, where $s$ represents the desired extension ratio for the context window.

The parameter $\hat{n}$ determines how many of the first token positions preserve the original RoPE embedding and are exempted from position interpolation. In practice, $\hat{n}$ is chosen from the set \{0, 1, 2, 4, 8, 12, 16, 20, 24, 28, 32, 64, 128, 256\}. If $\hat{n}=0$, then rescale factors obtained through the search process are applied to every token position.

\textbf{Evolution-based search}. The search space detailed in Table~\ref{tbl:searchspace} encompasses a vast array of positional interpolation configurations, making exhaustive exploration highly impractical. For instance, when the extension factor is $s=4\times$, the number of possible combinations reaches $400^{128/2}\times14$=4$\times10^{167}$, highlighting the immense complexity involved. This complexity grows exponentially as the extension ratio increases. To efficiently navigate this search landscape, we adopt an evolution search strategy~\cite{guo2020single} and propose two optimization methods that substantially enhance search performance. The comprehensive procedure is outlined in Algorithm~\ref{alg:search}.

\textit{Optimized initial population generation}. Rather than assigning all $P$ rescale factors at random, our approach incorporates the three RoPE rescale factors for PI, NTK, and YaRN directly as individuals within the starting population. The rest of the initial population, comprising $P-3$ individuals, is formed by introducing random mutations to these three rescale factors, each with a probability $p$.

\textit{Monotonically non-decreasing constraint}. Once the initial population has been created, we evaluate the LLM perplexity for every member. This involves adjusting the target LLM with the specified RoPE rescale factors and calculating the perplexity of the input $\mathbf{X}$. The best $k$ individuals, ranked by their perplexity scores, are selected as parents for the next generation. Nonetheless, due to the immense size of the search space, straightforward approaches to mutation and crossover may yield suboptimal candidates, resulting in excessive perplexity evaluations. This problem becomes especially pronounced when $L'$ is large, since each perplexity evaluation is computationally intensive.

To tackle this issue, we enforce a constraint of non-decreasing monotonicity on the sampled RoPE rescaling factors, such that $\lambda_i\le \lambda_{i+1}$. Only those RoPEs meeting this criterion are utilized when applying them to the LLM for perplexity computation, thereby substantially narrowing the search space. More concretely, we mandate that $\lambda_i$ progresses monotonically with the RoPE index (where $i=0,\ldots,63$). The rationale for this monotonicity across dimensions comes from NTK theory~\cite{jacot2018neural,tancik2020fourier,ntk}, which posits that dimensions with higher frequency (i.e., smaller $i$) should have reduced interpolation requirements (smaller $\lambda_i$), while those with lower frequency (larger $i$) can accommodate greater interpolation (larger $\lambda_i$).

\begin{algorithm}[t]
	\small
	\caption{The search algorithm}
	\textbf{Input:} target LLM, input samples $\mathbf{X}$, population size $P$, mutation size $N_1$, crossover size $N_2$, maximum iterations $\mathcal{T}$, mutation probability $p$\\

	\begin{algorithmic}[1]
		\label{alg:search}
		\STATE Top-k $\gets \phi$;	
		\STATE $\text{P}_0 \gets \textit{Initialize\_population\_with\_optimization}(P, \mathbf{X}, p)$; 
		\FOR{$i = 1$ \textbf{to} $\mathcal{T}$}
		\STATE \textit{Compute\_perplexity}(LLM, $\text{P}_{i-1}$, $\mathbf{X}$);
		\STATE Top-k $\gets$ \textit{Update\_Topk}(Top-k, $\text{P}_{i-1}$);
		\STATE $\text{P}_{mutation} \gets \textit{Mutation\_with\_mono\_constraint}$(Top-k, $N_1$, $p$); 
		\STATE $\text{P}_{crossover} \gets \textit{Crossover\_with\_mono\_constraint}$(Top-k, $N_2$);
		\STATE $\text{P}_i \gets \text{P}_{mutation} \cup \text{P}_{crossover} \cup$ Top-k;
		\ENDFOR
		\STATE Return the individual in Top-k with the minimum perplexity;
	\end{algorithmic}
\end{algorithm}

\textbf{8$\times$ extension without fine-tuning}. Leveraging evolutionary search, our approach discovers optimal non-uniform RoPE rescale factors, which selectively maintain important dimensions and positions, thereby reducing information loss from interpolation. As illustrated in Fig.\ref{fig:non-ft}, this enables us to expand LLaMA2's context window from 4k to 32k tokens without any need for fine-tuning. Conversely, prior techniques like PI, as well as non-uniform NTK and YaRN, experience significant increases in perplexity following just a 2$\times$ context expansion.

\subsection{Extending LLM Context Window to 2048K}
\label{sec:extendto2048k}

\textbf{Progressive extension to 2048k}. Here, we present our approach for expanding the context window of pre-trained LLMs from the standard 4k to beyond 2048k tokens. As previously shown, our non-uniform positional interpolation is capable of delivering an 8$\times$ increase in context length without the need for any fine-tuning. However, achieving even larger extensions—such as up to 512$\times$—necessitates fine-tuning. A possible strategy involves selecting appropriate RoPE rescaling factors corresponding to the desired 2048k window, followed by model fine-tuning. Nevertheless, this method encounters significant obstacles due to the considerable computational resources required for training. Additionally, we have found that effective fine-tuning becomes increasingly difficult when working with such extensive extension ratios (see Appendix for details).

Fortunately, LongRoPE proves to be beneficial for both the baseline and the length-extended versions of LLM after fine-tuning. In light of this, we propose an efficient and gradual approach that reaches the desired 2048k target using merely 1k fine-tuning iterations, all conducted with training sequences capped at 256k tokens.

$\Diamond$ \textit{Scaling pre-trained LLM to 256k using LongRoPE search}. Using LLaMA2 as a case study, we perform a search procedure aimed at achieving context window sizes of 128k and 256k. At this phase, the context expansion corresponds to ratios of 32$\times$ and 64$\times$, respectively.

$\Diamond$ \textit{Fine-tuning to 256k}. In this stage, we further train the pre-existing LLM so that it can handle a 256k context window. Initially, we fine-tune LLaMA2 over 400 iterations while utilizing the RoPE rescaling factors tailored for 128k. Subsequently, we switch the RoPE rescaling factors to the 256k variant on the resultant checkpoint and continue with an extra 600 fine-tuning steps. This approach demonstrates superior efficiency compared to commencing fine-tuning directly at 256k.

$\Diamond$ \textit{Scaling a fine-tuned extended LLM to 2048k via LongRoPE search}. In the final stage, we apply a subsequent search procedure to the LLM that has already been fine-tuned with a 256k-length context window. This approach enables us to achieve an extraordinarily large context window of 2048k tokens, all without requiring additional rounds of fine-tuning. Consequently, the total extension factor reaches $512\times$.

\textbf{Restoration of performance in shorter context windows}. When expanding the context window up to 2048k, we observe a decline in effectiveness within the model’s initial context window range. This decrease in performance is a recognized drawback of positional interpolation~\cite{pi}, where mapping positions in higher dimensions into a compressed space inside the original window leads to suboptimal utilization of the position embeddings, impairing the model's capabilities. Specifically, extending by a ratio of 512$\times$ results in considerable congestion of position encodings inside the primary 4k context window.

To address this issue, we carry out an additional evolutionary search on the expanded LLM aimed at tuning the RoPE rescale factors specifically for shorter context lengths, such as 4k and 8k tokens. The maximum permitted value for the searched $\lambda$ is decreased, since shorter lengths necessitate less positional interpolation. At inference time, the LLM adapts the RoPE rescale factors accordingly in a dynamic manner.

\section{LongRoPE}

\label{sec:method}
Building on these insights, we propose LongRoPE. This approach initially develops an efficient search strategy designed to leverage the two forms of non-uniformity. Subsequently, this mechanism is employed to scale the LLM context window to lengths surpassing 2 million tokens.

\subsection{Problem Formulation}

The presence of these two forms of non-uniformity considerably expands the solution space, thereby complicating the optimization process. To tackle this challenge, we recast the multidimensional non-uniform position interpolation optimization task as a search problem.

Given a large language model with a target context window length of $L'$ and input documents $\mathbf{X}$ composed of sequences $\mathbf{x}\in\mathbf{X}$ whose token counts exceed $L'$, we represent the original rotary position encoding angle for the $i^{th}$ dimension at token index $n$ as $\frac{n}{\beta^i}$. Accordingly, the optimization problem can be formally represented as follows:

\begin{equation}
\begin{gathered}
		\label{eq:problem}

\underset {\mathbf{x}\in\mathbf{X}; \, |\mathbf{x}|\ge L'}{ \operatorname {arg\,min} } \,  \mathcal{L}\, \left( \text{LLM} (\text{RoPE}, \mathbf{X})\right), \quad \text{where }
\\
\scalebox{0.93}{
$\underset { \begin{subarray}{l} i=0,\cdot\cdot,\frac{d}{2}-1; \\ n\in[0,|\mathbf{x}|);\end{subarray} }{\text{RoPE}_ {(n)} } = \left[ \cdots, \cos \left( \mathbb{I}(\hat{\lambda}_i, \hat{n}) \cdot \frac{n}{\beta^i} \right), \; \sin \left( \mathbb{I}(\hat{\lambda}_i, \hat{n}) \cdot \frac{n}{\beta^i} \right), \cdots \right] $ }
\\
\text{with } \mathbb{I}(\hat{\lambda}_{i},\hat{n}) =
\begin{cases}
    1 & \text{if } n < \hat{n} \\
    \frac{1}{\lambda_i} & \text{if } n \geq \hat{n}
\end{cases}
\end{gathered}
\end{equation}

In this approach, we incorporate a collection of rescale factors, $\mathbb{I}(\hat{\lambda}_{i},\hat{n})$, to address both types of non-uniformity. Here, $\hat{\lambda}_i$ represents non-uniform scaling across RoPE dimensions, while $\hat{n}$ accounts for non-uniformity in token positions. More precisely, $\mathbb{I}(\hat{\lambda}_{i},\hat{n})$ is employed to adjust the rotation angle in the $i^{th}$ RoPE dimension by utilizing $\hat\lambda_{i}$ as the scaling parameter and $\hat{n}$ as the critical token position boundary. For token positions less than $\hat{n}$, i.e., the first $\hat{n}-1$ tokens, the rotation angle remains unchanged as $\frac{n}{\beta^i}$, and no rescaling is applied. Once the token position reaches $n \ge \hat{n}$, the rescale factor $\hat{\lambda}_i$ is introduced.

With a specified target context window of $L'$, our goal is to determine the set of optimal rescaling coefficients ($\mathbb{I}(\hat{\lambda}_{0},\hat{n})$, $\mathbb{I}(\hat{\lambda}_{1},\hat{n})$, ..., $\mathbb{I}(\hat{\lambda}_{i},\hat{n})$, etc.) across the RoPE dimensions from the $1^{st}$ to the $d^{th}$. This optimization aims to ensure that the $\text{LLM}$, after applying these adjusted $\text{RoPE}$ scaling factors, attains the lowest possible next-token prediction loss, denoted by $\mathcal{L}$ (equivalent to minimal perplexity), on input sequences $\mathbf{X}$ that consist of $L'$ tokens.

\subsection{Searching the Non-uniform Position Interpolation}

To address the challenge posed in Eq.~\ref{eq:problem}, we present an efficient and straightforward approach that seeks optimal RoPE rescale factors, thereby maximizing the utilization of multidimensional irregularities present in position embeddings.

\textbf{Search space}. We construct an extensive search space that comprehensively incorporates the two distinct non-uniformities, as detailed in Table~\ref{tbl:searchspace}. More precisely, our method involves identifying a unique rescale factor for each individual dimension of the RoPE embedding. In order to simplify the formulation of the search space, our search is performed over $\lambda_i$ and $\hat{n}$ instead of directly searching for $\mathbb{I}(\hat{\lambda}_{i},\hat{n})$, with  $\hat{\lambda}_{i}=1/\lambda_i$. 
According to Table~\ref{tbl:searchspace}, the domain for $\lambda_i$ begins at 1.0, corresponding to straightforward extrapolation, and extends up to $s\times1.25$, which enables broader interpolation compared to PI, progressing in increments of 0.01. Here, $s$ represents the intended ratio for extending the context window.

The parameter $\hat{n}$ determines how many of the early token positions preserve the original RoPE embedding, foregoing position interpolation. In practice, $\hat{n}$ is selected from the set \{0, 1, 2, 4, 8, 12, 16, 20, 24, 28, 32, 64, 128, 256\} during the search process. Notably, if $\hat{n}=0$, every token position is assigned rescale factors determined by the search.

\textbf{Evolution-based search}. The search space outlined in Table~\ref{tbl:searchspace} includes a vast array of positional interpolation options, making it highly challenging to explore efficiently. For instance, an extension ratio of $s=4\times$ results in $400^{128/2}\times14$ = 4$\times10^{167}$ possible configurations. As the extension ratio increases, the size of the search space grows exponentially. To cope with this complexity, we leverage evolution search~\cite{guo2020single} and incorporate two optimization strategies designed to substantially improve search efficiency. The complete search framework is detailed in Algorithm~\ref{alg:search}.

\textit{Optimized initial population generation}. Rather than selecting all $P$ rescale factors at random to form the initial population, we deliberately include the three RoPE rescale factors associated with PI, NTK, and YaRN as part of the initial individuals. The other $P$-3 individuals are produced by applying random mutations to these three rescale factors, each with a mutation probability $p$.

\textit{Monotonically non-decreasing constraint}. Once the initial population has been generated, we evaluate the LLM perplexity for each candidate. This involves applying the respective RoPE rescale factors within the target LLM and measuring the perplexity of the input $\mathbf{X}$. The highest-performing $k$ individuals are then selected as parents for the subsequent evolutionary step. Nonetheless, due to the vastness of the solution space, simple mutation and crossover operations may frequently produce suboptimal candidates, incurring many redundant perplexity evaluations. Such inefficiency becomes especially pronounced when $L'$ is large, as each perplexity computation is computationally expensive.

To mitigate this, we enforce a constraint of non-decreasing monotonicity on the sampled RoPE rescaling factors such that $\lambda_i \leq \lambda_{i+1}$. Only those RoPE configurations adhering to this monotonic constraint are utilized when evaluating LLM perplexity, thereby substantially narrowing the search space and reducing computational expense. Concretely, our approach stipulates that the sequence $\lambda_i$ is monotonically increasing across the RoPE dimensions, for $i=0,\ldots,63$. This dimensional monotonicity is motivated by Neural Tangent Kernel (NTK) theory~\cite{jacot2018neural,tancik2020fourier,ntk}, which indicates that lower-indexed dimensions—characterized by higher frequency—require smaller scaling factors ($\lambda_i$) to minimize interpolation, while higher-indexed (lower-frequency) dimensions permit larger scaling factors, thus allowing for greater interpolation.

\begin{algorithm}[t]
	\small
	\caption{The search algorithm}
	\textbf{Input:} target LLM, input samples $\mathbf{X}$, population size $P$, mutation size $N_1$, crossover size $N_2$, maximum iterations $\mathcal{T}$, mutation probability $p$\\

	\begin{algorithmic}[1]
		\label{alg:search}
		\STATE Initialize Top-k as $\phi$;
		\STATE Set initial population $\text{P}_0$ via \textit{Initialize\_population\_with\_optimization} ($P$, $\mathbf{X}$, $p$); 
		\FOR{$i=1$ to $\mathcal{T}$}
		\STATE Execute \textit{Compute\_perplexity} (LLM, $\text{P}_{i-1}$, $\mathbf{X}$);
		\STATE  Update Top-k by invoking \textit{Update\_Topk} (Top-k, $\text{P}_{i-1}$);
		\STATE Derive $\text{P}_{mutation}$ using \textit{Mutation\_with\_mono\_constraint} (Top-k, $N_1$, $p$); 
		\STATE Generate $\text{P}_{crossover}$ using \textit{Crossover\_with\_mono\_constraint} (Top-k, $N_2$);
		\STATE Form next population $\text{P}_i$ by combining $\text{P}_{mutation}$, $\text{P}_{crossover}$, and Top-k;
		\ENDFOR
		\STATE Output the member of Top-k exhibiting the minimal perplexity;
	\end{algorithmic}
\end{algorithm}

\textbf{8$\times$ extension without fine-tuning}. By leveraging our evolutionary search strategy, we efficiently determine optimal non-uniform RoPE rescale factors, thereby maintaining critical dimensions and positional information to reduce information loss from interpolation. As illustrated in Fig.\ref{fig:non-ft}, our approach successfully extends the context window of LLaMA2 from 4k up to 32k tokens without requiring fine-tuning. Conversely, alternative techniques like PI, as well as non-uniform NTK and YaRN, result in significant increases in perplexity after merely doubling the context window size.

\subsection{Extending LLM Context Window to 2048K}
\label{sec:extendto2048k}

\textbf{Progressive extension to 2048k}. In this section, we present our approach for scaling the context window of existing pre-trained LLMs from the standard 4k up to beyond 2048k. As shown earlier, our non-uniform positional interpolation can yield up to an 8$\times$ increase without any additional fine-tuning. However, achieving even greater context lengths—such as a 512$\times$ extension—necessitates fine-tuning. One potential strategy involves identifying optimal RoPE rescaling factors for the desired 2048k window prior to commencing fine-tuning. Yet, this procedure is hindered by the immense computational resources required for training. Additionally, our practical experience indicates that effectively fine-tuning LLMs for such significant extension ratios is quite difficult (refer to Appendix).

Fortunately, LongRoPE demonstrates efficacy with both the base model and the LLM extended via fine-tuning. Accordingly, we propose a progressive and efficient approach that reaches the desired 2048k sequence length in only 1,000 fine-tuning iterations, all performed with a maximum training length of 256k.

$\Diamond$ \textit{LongRoPE search for extending pre-trained LLM to 256k}. Using LLaMA2 for illustration, we carry out a search aiming at context window sizes of 128k and 256k tokens. The resulting extension factors at these points are 32$\times$ and 64$\times$, respectively.

$\Diamond$ \textit{Fine-tuning to 256k}. In the subsequent stage, we adapt the pre-trained LLM to support a context window size of 256k. The process begins by fine-tuning LLaMA2 for 400 iterations with RoPE rescaling factors set for 128k. Following this, the RoPE rescaling factors are switched to accommodate 256k, after which a further 600 fine-tuning steps are performed starting from the previous checkpoint. This staged approach demonstrates greater efficiency compared to immediately fine-tuning for 256k.

$\Diamond$ \textit{Scaling up fine-tuned extended LLM to 2048k using LongRoPE search}. In the concluding phase, we conduct an additional search step with the LLM that was fine-tuned on a 256k sequence length. Through this process, we are able to enlarge the context window substantially, reaching 2048k, all achieved without the need for any more fine-tuning. This represents a context window extension factor of $512\times$.

\textbf{Recovery for shorter context windows}. When expanding the context window to a substantially large 2048k tokens, we observe that the model's effectiveness diminishes for inputs of the original window length. This outcome aligns with the documented challenges of positional interpolation~\cite{pi}, where embedding positional information into a higher-dimensional space causes the embeddings corresponding to positions in the initial context window to be compressed into a limited range, which impairs language modeling capabilities. Specifically, with an extension ratio of 512$\times$, position representations within the initial 4k context window become densely clustered.

To address this issue, we conduct an additional evolution search on the extended LLM specifically to calibrate the RoPE rescale factors for shorter context lengths (such as 4k and 8k). The upper bound for the searched $\lambda$ is lowered, as shorter contexts necessitate less positional interpolation. When inferring, the LLM adapts the relevant RoPE rescale factors in real time.

 \section{Experiments}

\label{sec:exp}
\subsection{Setup}
\textbf{Evaluation Tasks and Models}. Our experiments utilize LongRoPE with both LLaMA2-7B and Mistral-7B, assessing their capabilities across three dimensions: (1) the perplexity of large language models with extended context lengths when processing lengthy texts; (2) a Passkey retrieval task designed to assess how effectively a model can extract a specific passkey embedded among distractor passages; and (3) performance on conventional LLM benchmarks constrained to a 4096-token context window.

\textbf{Fine-tuning}. For LLaMA2, the fine-tuning procedure employs a learning rate of 2e-5 with a linear decay schedule and utilizes a global batch size of 32. The model is first trained for 400 steps on the Redpajama~\cite{redpajama} dataset, which is partitioned into 128k-token segments, each delimited by BOS and EOS tokens. Subsequently, using the resulting checkpoint, an additional 600 training steps are performed to extend the context window to 256k tokens. Training with 128k context is conducted using 8 A100 GPUs supported by the distributed training system~\cite{cube}, while expanding to 256k tokens requires scaling up to 16 A100 GPUs. For Mistral, a steady learning rate of 1e-6 and a global batch size of 64 are employed. Both the 128k and 256k context window variants follow the protocol described in YaRN~\cite{yarn}, involving 400 training steps on the Together Computer's Long-Data Collections~\cite{mistraldata} with a sequence length set to 16k tokens. The training for Mistral is performed on 4 A100 GPUs.

\textbf{Search}. For scenarios where the target window size does not exceed 256k, the following parameters are employed: $P$=64, $N_1$=$N_2$=16, $p$=0.3, and $\mathcal{T}$=40. In every iteration, the 32 highest-ranking candidates are chosen for mutation and crossover. Perplexity evaluation is performed on 5 randomly selected samples from the PG19 validation set, each meeting the minimum required context length. When the window size surpasses 512k, the values for population, mutation, and crossover are reduced by half. In this larger window regime, perplexity is assessed on 3 randomly chosen samples from the Pile-Books3~\cite{pile} validation split.

\textbf{Baselines}. To achieve a context window of 2048k, we conducted fine-tuning on models with 128k and 256k token contexts. These procedures resulted in the creation of LongRoPE-2048k (ft=128k) and LongRoPE-2048k (ft=256k) for LLaMA2 and Mistral, respectively. The performance of these four models was evaluated against leading baselines for expanding context windows in open-source LLMs, focusing on models fine-tuned after positional interpolation strategies such as PI, NTK, and YaRN. The baseline comparisons encompass models such as Together-32k~\cite{together}, Code LLaMA~\cite{codellama}, LongLoRA-full-FT-100k~\cite{longlora}, as well as YaRN-LLaMA and YaRN-Mistral~\cite{yarn}.

\subsection{Main Results}

\label{sec:mainresult}
\textbf{Long sequence language modeling up to 256k.} We initiate our analysis by benchmarking against leading extended LLMs at a context window of 256k tokens. To showcase the breadth of our method, evaluations are conducted on the Proof-pile~\cite{pg19} and PG19~\cite{pile} test splits. Perplexity is measured across different context lengths utilizing a sliding window of 256 tokens. For PG19, the complete test split comprising 100 documents is employed. For Proof-pile, we adopt the sampling methodology from YaRN~\cite{yarn}, selecting 10 samples at random, each exceeding 128k tokens in length.

Table~\ref{tbl:proof-pile} and Table~\ref{tbl:pg19} present perplexity comparisons between LLaMA2 and Mistral models augmented with various interpolation strategies, evaluated on Proof-pile and PG19, respectively. Two main findings emerge: \textbf{(1)} our enhanced models consistently exhibit reduced perplexity as evaluation lengths increase from 4k up to 256k, indicating that they effectively utilize longer contexts. \textbf{(2)} Despite operating within an expanded context window—specifically, 16$\times$ longer than the standard setting, which generally complicates maintaining short-range performance—our LongRoPE-2048k variants still surpass leading baseline models for context lengths up to 256k.

\textbf{Long sequence language modeling beyond 2000k}. To assess performance on exceptionally lengthy documents, we employ the Books3~\cite{pile} dataset. For efficient evaluation, we randomly sample 20 books, each with a length greater than 2048k, and apply a 256k sliding window.

As presented in Table~\ref{tbl:books3}, LongRoPE is able to expand the context windows of LLaMA2-7B and Mistral-7B models up to 2048k, achieving perplexity scores that are on par with, or even better than, those of baseline methods at shorter sequence lengths ranging from 8k to 128k. A notable difference is observed between the LLaMA2 and Mistral models at the 2048k context length. Specifically, Mistral surpasses baseline models for shorter contexts; however, its perplexity rises above 7 for sequence lengths exceeding 256k. In the case of LLaMA2, model performance follows anticipated trends—perplexity decreases steadily as the context lengthens, with only slight increases at 1024k and 2048k. Additionally, for LLaMA2, LongRoPE-2048k achieves better results when fine-tuned at a length of 256k versus 128k, which is attributed to the smaller extension multiplier (8$\times$ compared to 16$\times$). Conversely, Mistral attains superior performance when fine-tuned with a window size of 128k. This outcome primarily stems from the fact that for Mistral's fine-tuning at both 128k and 256k, a 16k training length (as recommended by YaRN) is employed, thereby limiting Mistral’s potential for context window extension post fine-tuning.

\textbf{Passkey retrieval}. Next, we investigate the functional context window length in generative scenarios. To do so, we adopt the synthetic passkey retrieval task introduced by~\cite{passkey}. In this setup, the language model's goal is to extract a randomly selected five-digit passkey embedded within a lengthy document. A full description of the prompt template is provided in the appendix. For the evaluation, we conduct 10 rounds of passkey retrieval, positioning the passkey at uniformly random locations throughout the span of the evaluation context in each iteration.

Fig.~\ref{fig:passkey} presents a comparison of retrieval accuracy across different models and baselines. The accuracy of previous approaches falls sharply to zero once the input surpasses 128k tokens. By contrast, even when tasked with retrieving a passkey from millions of tokens—a significantly more difficult scenario—LongRoPE-LLaMA2-2048k (ft=256k) consistently delivers high retrieval accuracy ($\ge$90\%) across the range from 4k to 2048k. Similarly, LongRoPE-Mistral-2048k (ft=128k) sustains perfect (100\%) accuracy up to 1800k tokens before decreasing to 60\% at 2048k, which is in line with the increase in perplexity observed at 2048k, as reported in Table~\ref{tbl:books3}.

\textbf{Evaluation on standard benchmarks within the original context window}. LongRoPE-2048k models are assessed on the original context window employing the Hugging Face Open LLM Leaderboard~\cite{huggingfaceboard} in both zero-shot and few-shot paradigms. Specifically, we adopt 25-shot ARC-Challenge~\cite{allenai:arc}, 10-shot HellaSwag~\cite{zellers2019hellaswag}, 5-shot MMLU~\cite{mmlu}, and 0-shot TruthfulQA~\cite{lin2021truthfulqa} tasks for this evaluation.

As demonstrated in Table~\ref{tbl:commonsense}, our models perform at a similar level to the original benchmark, which was created for a more limited context window, and surpass the original Mistral on TruthfulQA by 0.5\%. While LongRoPE-LLaMA2-2048k, which was fine-tuned using 256k tokens, exhibits a minor drop in performance, it still maintains acceptable results across the majority of evaluated tasks.

\subsection{Ablation Results}

\textbf{Effectiveness of the second positional interpolation}. Within our progressive extension framework, we implement our search algorithm to perform a secondary non-uniform positional interpolation on already fine-tuned and expanded LLMs. To assess its efficacy, we carry out experiments using the LLaMA2-256k model after fine-tuning. The model is subsequently expanded to 512k, 1024k, and 2048k contexts through both PI and YaRN methods. As presented in Table~\ref{tbl:secondsearch}, our approach with non-uniform positional interpolation maintains stable perplexity scores throughout the extension process. By comparison, applying PI and YaRN results in rapidly worsening perplexity as the extension ratio increases.

\textbf{Effectiveness of recovery at shorter context lengths.} To address reductions in performance when handling shorter context lengths, we recalibrate the RoPE factors for LongRoPE-2048k using our search-based optimization procedure. In particular, we restrict the upper bound of scale factors in the search space, thereby reducing the degree of interpolation at the 4k and 8k context lengths. Table~\ref{tbl:shortperformance} presents a comparison of perplexity for LongRoPE-LLaMA2-2048k evaluated on Proof-pile at both 4k and 8k, together with the average LLM benchmark accuracy. The data indicates a substantial enhancement in model performance at these shorter context sizes.

\textbf{Analysis on the two forms of non-uniformities}. To disentangle the effects of the two non-uniformities, we conduct ablation studies to examine their respective contributions to overall performance. Two sets of experiments are designed: (i) upscaling LLaMA2-7B to accommodate sequence lengths of 16k and 32k using different strategies—PI, optimizing solely the RoPE dimension, and jointly optimizing both non-uniformities; (ii) expanding the input length of a fine-tuned 256k LLaMA2 to 2048k, employing analogous procedures. For both scenarios, perplexity is measured without further fine-tuning. According to the results in Table~\ref{tbl:searchdimension}, introducing non-uniformity within the RoPE dimension yields a notable reduction in perplexity compared to simple linear interpolation with PI. Meanwhile, non-uniform token position confers distinct gains for shorter sequence lengths of 16k and 32k, though similar benefits do not appear at the extreme length of 2048k—potentially because retaining the initial set of tokens without performing interpolation is ineffective in such scenarios. Further exploration of these effects is proposed as future work.

\section{Related Works}

In addition to methods based on position interpolation, this section discusses related works of other approaches. 

\textbf{Retrieval-based} methods employ external memory components to store distant contextual information and utilize retrieval mechanisms to access pertinent documents during inference~\cite{longllama,wang2023augmenting,borgeaud2022improving}. Such strategies often require substantial and explicit changes to the underlying LLM architectures. In comparison, our approach remains relatively lightweight, introducing only slight alterations to the positional embedding scheme. Furthermore, our method is applicable to a wider range of long context tasks, extending beyond retrieval to include scenarios like long document summarization and few-shot learning.

\textbf{Attention-based context window extensions}. In addition to positional embedding interpolation, alternative methods have extended the effective input context by adjusting attention mechanisms within the fixed context window size of the original LLM~\cite{han2023lminfinite, streamingllm,parallelwindow}. Central to these approaches is the use of innovative attention masks designed to address the increase in attention computation associated with added positions. These strategies can be combined with positional interpolation techniques to provide complementary benefits.

\textbf{Fine-tuning based approaches} are primarily concerned with strategies for adapting pre-trained LLMs to accommodate longer sequences, typically by adjusting position embeddings. Prior work, such as Code LLaMA~\cite{codellama}, LLaMA2 Long~\cite{xiong2023effective}, and ScaledRoPE~\cite{liu2023scaling}, utilize a significantly increased base value for RoPE, followed by fine-tuning at the required sequence length. In contrast, our approach provides adaptability across a range of target lengths and demonstrates effectiveness at lengths surpassing 2M tokens. More recently, to address the substantial computational costs associated with fine-tuning on very long contexts (greater than 128k), methods such as LongLoRA~\cite{longlora} and PoSE~\cite{zhu2023pose} have been introduced to reduce this resource demand. Our proposed technique is complementary to these efficient fine-tuning methods.

\section{Conclusion}

This paper introduces LongRoPE, a technique that dramatically increases the context length of LLMs up to 2048k, while preserving performance on the original, shorter context spans. The approach harnesses two types of non-uniformities inherent in RoPE positional embeddings, identified via a highly efficient evolutionary search algorithm. This methodology yields two major advantages: first, it allows for robust initialization conducive to fine-tuning, and second, it permits extending the context window by a factor of 8 without any fine-tuning required. Furthermore, a staged extension approach is put forward, wherein LLMs fine-tuned for a 256k context length are progressively adapted to operate at 2048k without needing further fine-tuning. Comprehensive experimental results underscore LongRoPE’s efficacy. We anticipate that our LongRoPE-2048k models will unlock a range of new applications requiring extended context windows and spur continued innovation in this domain.

\section*{Broader Impacts} The research detailed in this paper aims to contribute to the ongoing progress in Machine Learning as a discipline. While our findings may have a range of societal implications, we do not identify any particular issues requiring emphasis at this time.

	\balance

\end{document}